\bookchapter{บทที่ 67: ขัดเกลาและเผยแพร่ — ทดสอบ ดีบัก และขึ้น App Store}

\chapterintro{ทำแอปเสร็จแล้ว แล้วต่อไปล่ะ? บทนี้ตอบสามเรื่องของการไปจาก'รันได้'สู่'ส่งมอบได้': ดีบักหาบั๊กอย่างมีประสิทธิภาพอย่างไร เขียนเทสต์กันบั๊กย้อนกลับอย่างไร และส่งแอปขึ้น App Store ให้ทั่วโลกดาวน์โหลดอย่างไร การขึ้นร้านไม่ใช่แค่พิธีกรรม — มันคือเส้นแบ่งที่เปลี่ยนผลงานชิ้นหนึ่งให้เป็นผลิตภัณฑ์}

\booksection{สามท่าเด็ดการดีบัก: print, จุดพัก, พรีวิว}

\begin{center}
\begin{tabularx}{\linewidth}{|X|X|X|}
\hline
\textbf{เครื่องมือ} & \textbf{ทำหน้าที่อะไร} & \textbf{ใช้อย่างไร} \tabularnewline
\hline
print + คอนโซล & ดูล็อกการทำงาน & ในโค้ดพิมพ์ print("ค่า=\textbackslash{}(x)") แล้วดูผลลัพธ์ที่คอนโซลด้านล่างของ Xcode \tabularnewline
การดีบักด้วยจุดพัก & ให้โปรแกรมหยุดที่บรรทัดใดบรรทัดหนึ่งเพื่อดูค่าตัวแปร & คลิกข้างเลขบรรทัดจะมีตัวหมายสีน้ำเงิน รันถึงจุดนั้นจะหยุด แล้วไล่ทีละขั้น \tabularnewline
พรีวิว SwiftUI + ดีบักเกอร์ & ดูอินเทอร์เฟซและลำดับชั้นวิวแบบเรียลไทม์ & แคนวาส \#Preview + ตัวตรวจ View Hierarchy ตอนรันจริง \tabularnewline
\hline
\end{tabularx}
\end{center}

\begin{center}\small\color{gray}สามท่าเด็ดการดีบัก\end{center}

\begin{tcbinfo}{เคล็ดลับ: แอปพังแล้วอ่านอย่างไร?}
ตอน Swift พัง คอนโซลจะพรินต์ “Fatal error” พร้อมสาเหตุ และ Xcode จะหยุดที่บรรทัดที่มีปัญหาพร้อมไฮไลต์ไว้ ประวัติการพังบนเครื่องจริงดูได้ที่ Xcode Window Organizer Crashes (ต้องใช้บัญชีแบบเสียเงิน) ตอนค้นหาให้เอา'ประโยคแรกของสาเหตุ error + swiftui'ไปค้นก่อน — คำถามแม่น คำตอบถึงจะแม่น

\end{tcbinfo}

\booksection{การทดสอบ: ทำประกันให้โค้ด}

การเทสต์ด้วยมือไม่น่าเชื่อถือ — แก้ฟีเจอร์นี้ไป อีกตัวเงียบ ๆ พังไป เทสต์จึงเป็นทั้ง'ยาแก้สำนึก'และ'เสื้อกันกระสุน' iOS ใช้เฟรมเวิร์ก XCTest แบ่งเป็นสองประเภท:

\begin{center}
\begin{tabularx}{\linewidth}{|X|X|X|X|}
\hline
\textbf{ระดับ} & \textbf{ทดสอบอะไร} & \textbf{Target} & \textbf{ความเร็ว} \tabularnewline
\hline
Unit Test & โลจิกล้วน ๆ (การคำนวณ, การตรวจสอบ) & MyFirstAppTests & ระดับวินาที \tabularnewline
UI Test & พฤติกรรมอินเทอร์เฟซจริง & MyFirstAppUITests & ระดับนาที \tabularnewline
เทสต์ด้วยมือ & ความรู้สึกและประสบการณ์ของคนจริง & —— & ต้องมีตลอดไป \tabularnewline
\hline
\end{tabularx}
\end{center}

\begin{center}\small\color{gray}ระดับการทดสอบบน iOS (โปรเจกต์เทมเพลตมี Test Target ให้สองตัวแล้ว)\end{center}

\begin{lstlisting}[language=Swift]
import XCTest

// แยกตรรกะ "แปลงข้อความจำนวนเงินเป็นตัวเลข" ออกเป็นฟังก์ชันล้วน
func parseAmount(_ input: String) -> Double? {
    guard let v = Double(input.trimmingCharacters(in: .whitespaces)),
          v > 0 else { return nil }
    return v
}

final class AmountTests: XCTestCase {
    func testEmptyInputReturnsNil() {
        XCTAssertNil(parseAmount(""))
    }
    func testNegativeReturnsNil() {
        XCTAssertNil(parseAmount("-5"))
    }
    func testNormalValueParses() {
        XCTAssertEqual(parseAmount("12.5"), 12.5)
    }
}
\end{lstlisting}

รันเทสต์: CmdU ให้เขียวหมดก่อนแล้วค่อยแก้โค้ด ใจก็นิ่ง การตรวจจำนวนเงินและการคำนวณสรุปยอดในสมุดบัญชี ควรมี Unit Test คุ้มกันไว้

\booksection{การเซ็นและแพ็กเกจ: ประทับ'ตรา'ให้แอป}

แอป iOS ทุกตัวต้องเซ็นก่อนถึงจะติดตั้งได้ — เพื่อพิสูจน์ว่า'แพ็กเกจนี้มาจากคุณและไม่ถูกแก้ไข' ความสัมพันธ์แบบลำดับชั้น:

\begin{itemize}
  \item ช่วงเรียน: Apple ID ฟรี + Xcode เซ็นให้อัตโนมัติ (Personal Team) รันบนเครื่องจริงได้ ต้องต่ออายุทุก 7 วัน
  \item ช่วงเผยแพร่: ต้องใช้บัญชีนักพัฒนาแบบเสียเงิน (99 ดอลลาร์/ปี) ออกใบรับรองและโปรไฟล์จัดเตรียม จึงจะขึ้นร้านและติดตั้งระยะยาวได้
\end{itemize}

แพ็กเกจ: เมนู Xcode Product Archive (เก็บถาวร) พอบิลด์เสร็จไปที่ Organizer เลือก Distribute App เลือก App Store Connect อัปโหลด การเก็บถาวรต้องตั้งเป้าหมายเป็น Generic iOS Device (ไม่ใช่โปรแกรมจำลอง) นี่คือขั้นตอนที่มือใหม่ลืมบ่อยที่สุด

\textbf{กันพลาด: ใบรับรองหาย = การอัปเดตสะดุด}\par

ใบรับรองเซ็นและคีย์ส่วนตัวพอหายไป คุณจะส่งอัปเดตให้แอปเก่าไม่ได้อีก — ผู้ใช้จะติดตั้งได้แต่แอป'ใหม่เอี่ยม' สำรองใบรับรองให้ดี (ส่งออก .p12 จากคีย์เชน) และจดรหัสผ่านไว้ให้แน่น โชคดีที่ Apple มี'การเซ็นอัตโนมัติ'และการจัดการใบรับรองผ่านคลาวด์ (ติ๊กเลือก Automatically manage signing ใน Xcode) ช่วยตัดปัญหางานมือไปได้ส่วนใหญ่

\booksection{ขั้นตอนครบวงจรการขึ้น App Store}

สมัครบัญชีนักพัฒนา

เตรียมสื่อสำหรับสโตร์

อัปโหลดแพ็กเกจบิลด์

กรอกข้อมูลความเป็นส่วนตัวและการรีวิว

รีวิวโดย Apple (1-7 วัน)

เผยแพร่บนสโตร์

\begin{center}\small\color{gray}ภาพที่ 5: ขึ้น App Store หกขั้นตอน\end{center}

\begin{itemize}
  \item ลงทะเบียน: developer.apple.com บัญชีบุคคล 99 ดอลลาร์/ปี
  \item สื่อประกอบ: ไอคอน (1024×1024 อย่าทำมุมโค้ง — Apple ตัดเอง), ภาพหน้าจอ (iPhone ทุกขนาด), คำอธิบาย, ลิงก์นโยบายความเป็นส่วนตัว
  \item แบบสอบถามความเป็นส่วนตัว: แจ้งตรง ๆ ว่าเก็บข้อมูลอะไรบ้าง — สมุดบัญชีที่ไม่ต่อเน็ตเลือก'ไม่เก็บข้อมูลใด ๆ'ได้
  \item การรีวิว: ส่งครั้งแรกมักต้องรอ 1-7 วัน ถูกปฏิเสธก็แก้ตาม'หมายเหตุจากทีมรีวิว'ทีละข้อ — เหตุผลการปฏิเสธของ Apple มักเขียนละเอียดมาก ทำตามก็จบ
  \item TestFlight: ก่อนขึ้นร้านส่งเวอร์ชันทดสอบให้เพื่อนทดลองใช้ก่อน (ได้สูงสุด 10000 คน) เก็บฟีดแบ็กหนึ่งรอบแล้วค่อยส่งจริง ดูมืออาชีพสุด ๆ
\end{itemize}

\booksection{ทำสื่อหน้าร้านอย่างไร: ความประทับใจแรกคือจำนวนดาวน์โหลด}

หน้าร้านคือช่องโชว์ที่ผู้ใช้เห็นตอน'เดินเลือกซื้อ' สื่อสามชิ้นกำหนดความประทับใจแรก:

\begin{itemize}
  \item ไอคอน: 1024×1024 อย่าใช้ภาพหน้าจอมาแกะโล่ะ ถ้าไม่มีฝีมือออกแบบ ใช้เครื่องมือ AI สร้างกราฟิกเรียบ ๆ สักชิ้น ขอให้ย่อลงบนหน้าจอโฮมแล้วยังมองปุ๊บรู้เรื่องทันทีถึงจะผ่าน
  \item ภาพหน้าจอ: อย่างน้อย 2 ชุด (ขนาดยอดนิยมของ iPhone) โชว์ฟีเจอร์หลักที่สำคัญที่สุด แต่ละภาพแนบคำโปรยภาษาคนหนึ่งประโยค ('จดหนึ่งรายการในหนึ่งวินาที'ดีกว่า'โซลูชันบันทึกบัญชีที่มีประสิทธิภาพ'สิบเท่า)
  \item คำอธิบาย: ประโยคแรกบอกให้ชัดว่า'นี่คืออะไร + ทำให้ใครใช้' อย่ากองคำว่า'มืออาชีพที่สุด เก่งที่สุด' — ความจริงใจได้ผลกว่าความเว่อร์
\end{itemize}

\begin{tcbpractice}{ลงมือทำ: 'ส่งมอบ'สมุดบัญชีสักหนึ่งครั้ง}
\booksection{รายการสั้น ๆ เรื่องการรายงานการพังและการปรับประสิทธิภาพ}

ก่อนแอปขึ้นร้าน สองเรื่องนี้คุ้มที่จะใช้เวลาคนละครึ่งชั่วโมง:

\booksubsection{การเฝ้าระวังการพังของแอป}

ตอนคุณทดสอบไม่พัง ไม่ได้แปลว่าในมือถือผู้ใช้จะไม่พัง Apple มีรายงานการพังฟรีใน Xcode Organizer (บัญชีแบบเสียเงินดูล็อกการพังบนเครื่องจริงได้) ทางที่สบายใจกว่าคือต่อ SDK ภายนอก (อย่าง Firebase Crashlytics) ขึ้นร้านโดยไม่มีการเฝ้าระวังการพัง เท่ากับขับรถลืมตาปิด

\booksubsection{ตรวจประสิทธิภาพเองสี่ข้อ}

\begin{enumerate}
  \item รายการใช้ LazyVStack / List: ลิสต์ยาว ๆ อย่าจัดด้วย VStack ธรรมดาเป็นอันขาด — การโหลดแบบขี้เกียจคือบทเรียนแรกของประสิทธิภาพ
  \item รูปอย่ายัดไฟล์ต้นฉบับใหญ่เข้าไป: ใช้ AsyncImage (รูปจากระยะไกล) กับแคชของระบบ ไฟล์ต้นฉบับ 5MB พอจะทำเครื่องรุ่นเก่าค้างทั้งตัว
  \item อย่าทำงานหนักบนเธรดหลัก: ทั้งเครือข่ายและไฟล์ใหญ่ให้ผ่าน async/await หมด
  \item ใช้ Instruments วัดดู: ตัววิเคราะห์ประสิทธิภาพที่มากับ Xcode (Time Profiler, ตรวจหาหน่วยความจำรั่ว) — วัดก่อน แล้วค่อยปรับ อย่าเชื่อความรู้สึก
\end{enumerate}

(1) เพิ่ม Unit Test ให้โลจิกอย่าง parseAmount สามตัว (รัน CmdU ให้เขียว); (2) ใช้จุดพักจับการพังที่ตัวเองสร้างขึ้นเองสักหนึ่งครั้งแล้วแก้ให้ดี; (3) Product Archive แพ็กเกจหนึ่งชุด ผ่าน TestFlight เชิญเพื่อนมาทดลองใช้; (4) ถ้ามีกำลังเหลือ ลงทะเบียนบัญชีนักพัฒนาแล้วส่งสมุดบัญชีขึ้นร้านจริง ๆ สามขั้นแรกทำเสร็จได้ภายในคืนนี้

\end{tcbpractice}

\begin{tcbsummary}{สรุปบทนี้}
\begin{itemize}
  \item สามท่าเด็ดการดีบัก: ล็อก print, จุดพัก, พรีวิว SwiftUI กับตัวตรวจวิว
  \item XCTest: Unit Test (CmdU) + UI Test + เทสต์ด้วยมือ เทมเพลตมีให้พร้อมแล้ว
  \item การเซ็น: บัญชีฟรีเซ็นให้อัตโนมัติ (ต่ออายุทุก 7 วัน); แจกจ่ายจริงต้องใช้บัญชีแบบเสียเงิน
  \item แพ็กเกจ: Generic iOS Device + Archive + Distribute; สำรองใบรับรองให้แน่นอน
  \item หกขั้นสู่ร้าน: ลงทะเบียน สื่อประกอบ อัปโหลด แบบสอบถามความเป็นส่วนตัว รีวิว วางขาย; ให้ TestFlight ทดลองใช้ก่อน
\end{itemize}

\end{tcbsummary}
